\documentclass[11pt,a4paper]{article}

\input{../preambule}

\begin{document}

\title{Semaine 2}
\author{}
\date{28/09 - 04/10 }

\maketitle


\begin{exercise}
Expliciter géométriquement tous les $z\in \C$ tels que $z, z^2, z^3$ forment un triangle isocèle.
\end{exercise}


\begin{exercise}
\begin{questions}
\item Soit $n\in \N^*$. Résoudre l'équation $(1+ix)^n=(1-ix)^n$ d'inconnue $x\in \R$.
\item En déduire la valeur de $\displaystyle\sum_{k=0}^{n-1} \tan\left(\frac{k\pi}{n}\right)$ si $n$ est impair.
\end{questions}
\end{exercise}




\end{document}